\chapter{Métodos Numéricos Auxiliares}
\label{ap:metodos-auxiliares}

\section{Ortogonalización de Gram-Schmidt}
\label{sec:gram-schmidt}

En un cálculo Hartree-Fock estándar para configuraciones de capas cerradas, los orbitales espaciales radiales $P_{nl}(r)$ que comparten el mismo momento angular $l$ pero difieren en el número cuántico principal $n$ emergen automáticamente ortogonales, es decir, $\langle P_{nl} | P_{n'l} \rangle = \delta_{nn'}$. Esto ocurre porque, matemáticamente, dichos orbitales son los autovectores correspondientes a autovalores distintos ($\epsilon_{nl} \neq \epsilon_{n'l}$) de un \textit{mismo} operador hermitiano (la matriz de Fock $\mathbf{F}$). El teorema espectral garantiza su ortogonalidad intrínseca.

Sin embargo, en el esquema de Hartree-Fock Restringido para Capas Abiertas (ROHF) que implementamos para abordar configuraciones como $np^2$, la simetría esférica del potencial colapsa. Para mantener la forma central de los orbitales, el operador de Fock se divide en Hamiltonianos efectivos disjuntos para las diferentes capas. Específicamente, en nuestro motor computacional, las capas cerradas (core) se optimizan bajo un operador $\mathbf{F}_c$, mientras que la capa de valencia abierta obedece a un operador distinto $\mathbf{F}_v$.

Al provenir de diagonalizaciones de operadores numéricos diferentes, se pierde la garantía espectral estricta: un orbital de valencia $P_{v}$ no será perfectamente ortogonal a un orbital del core $P_{c}$ con el mismo momento angular, provocando solapamientos espurios $S_{vc} = \langle P_v | P_c \rangle \neq 0$. Si este cruce no nulo entra al ciclo SCF, la matriz de densidad se contamina y el sistema diverge en espiral.

Para imponer coercitivamente el requerimiento físico del Principio de Exclusión de Pauli, implementamos un proceso de proyección en cada iteración, conocido como el método de \textbf{Ortogonalización de Gram-Schmidt}. Después de que el solucionador de eigenvalores obtiene el vector de la capa de valencia $| \tilde{P}_v \rangle$ expandido en la base B-spline, calculamos su superposición con todos los orbitales del core previamente fijados $| P_{c_i} \rangle$:

\begin{equation}
    | P_v \rangle = | \tilde{P}_v \rangle - \sum_{i \in \text{core}} \langle \tilde{P}_v | P_{c_i} \rangle | P_{c_i} \rangle.
\end{equation}

Computacionalmente, en el espacio discreto gobernado por la matriz de solapamiento $\mathbf{B}$, la operación se ejecuta mediante vectores columna expandidos:

\begin{equation}
    \mathbf{c}_v^{\text{ortogonal}} = \tilde{\mathbf{c}}_v - \sum_{i \in \text{core}} \left( \tilde{\mathbf{c}}_v^T \mathbf{B} \mathbf{c}_{c_i} \right) \mathbf{c}_{c_i}.
\end{equation}

Posterior a esta sustracción escalar, el orbital de valencia purificado es renormado $\mathbf{c}_v \leftarrow \mathbf{c}_v / \sqrt{\mathbf{c}_v^T \mathbf{B} \mathbf{c}_v}$ y reintegrado al ciclo SCF, asegurando convergencia incondicional \cite{Roothaan1960}.

\section{Cuadratura de Gauss-Legendre en B-splines}
\label{sec:gauss-legendre}

Los elementos matriciales de Galerkin $\int B_i(r)\,\mathcal{L}\,B_j(r)\,dr$ del Hamiltoniano y de la ecuación de Poisson se evalúan por cuadratura de Gauss-Legendre sobre cada intervalo de la malla \cite{Bachau2001}. Una regla gaussiana de $N_q$ puntos integra exactamente cualquier polinomio de grado hasta $2N_q - 1$. Como los B-splines de orden $k$ son polinomios por partes de grado $k-1$, los integrandos de las matrices de solapamiento y cinética son polinomios de grado a lo sumo $2k-2$ en cada intervalo, y el motor emplea $N_q = k+6$ puntos, de modo que esas dos matrices se obtienen sin error de cuadratura.

Los demás integrandos contienen factores $1/r$, $1/r^2$ o $r^{-4}$, o productos de tres funciones base con un factor $1/r$, y no son polinómicos. Para ellos la cuadratura no es exacta. Los seis puntos adicionales reducen su error, y la malla exponencial, que concentra intervalos cortos cerca del origen, donde esos factores varían más rápido, contribuye a mantenerlo pequeño.

Las reglas de cuadratura canónicas están definidas sobre el dominio estándar $[-1, 1]$ mediante la sumatoria ponderada:
\begin{equation}
    \int_{-1}^{1} f(x) dx \approx \sum_{q=1}^{N_q} w_q f(x_q),
\end{equation}
donde las raíces $x_q$ son los ceros del polinomio de Legendre $P_{N_q}(x)$ y $w_q$ son los pesos asociados.

Debido a que nuestros polinomios B-spline están definidos en los subintervalos nodales locales $[t_i, t_{i+1}]$ generados por nuestra malla no lineal discreta, procedemos a realizar un mapeo de transformación afín unidimensional. El cambio de variable para transferir la cuadratura del espacio estándar $x \in [-1, 1]$ al espacio físico $r \in [t_i, t_{i+1}]$ está dado por:

\begin{equation}
    r(x) = \frac{t_{i+1} - t_i}{2} x + \frac{t_{i+1} + t_i}{2}, \quad \text{con diferencial } dr = \frac{t_{i+1} - t_i}{2} dx.
\end{equation}

Con el jacobiano $J_i = \frac{t_{i+1} - t_i}{2}$, la integral sobre cada intervalo se aproxima por

\begin{equation}
    \int_{t_i}^{t_{i+1}} f(r) dr \approx J_i \sum_{q=1}^{N_q} w_q f(r(x_q)),
\end{equation}

con igualdad exacta cuando $f$ es un polinomio de grado a lo sumo $2N_q - 1$. Los valores de la base y de sus derivadas en los puntos $r(x_q)$ forman los tensores de colocación de la Sección \ref{sec:estructura-algebraica}.

\section{Reducción del Potencial de Polarización del Core al Caso Atómico}
\label{sec:cpp-muller}

El potencial de polarización empleado en el Capítulo \ref{ch:correlacion-estructura-fina} se presenta allí a partir de un argumento electrostático elemental: un electrón de valencia induce un dipolo en el core y se estabiliza al interactuar con él. Esa derivación es suficiente para justificar la forma $-\alpha_d/2r^4$, pero conviene mostrar que coincide con la formulación general del \textit{Core Polarization Potential} (CPP) de Müller, Flesch y Meyer \cite{Muller1984}, concebida para moléculas y agregados, cuando se le imponen las condiciones de un átomo aislado. El ejercicio deja explícito, además, un término de dos cuerpos que la derivación elemental no hace visible.

\subsection{La formulación general}

En su forma original, el potencial de polarización del sistema completo se escribe como una suma sobre centros polarizables $c$,
\begin{equation}
    \label{eq:cpp_general}
    V_{\text{CPP}} = -\frac{1}{2}\sum_c \left( \alpha_c\, \bm{f}_c \cdot \bm{f}_c - 2\,\bm{f}_c \cdot \bm{\mu}_c^{\circ} + \bm{f}_c^{\circ}\cdot\bm{\mu}_c^{\circ} \right),
\end{equation}
donde el campo eléctrico efectivo que actúa sobre el centro $c$ recibe contribuciones de todos los electrones y de todos los demás núcleos,
\begin{equation}
    \label{eq:cpp_campo}
    \bm{f}_c = \sum_i \frac{\bm{r}_{ci}}{r_{ci}^3}\,C(r_{ci}, \rho_c) - \sum_{c' \neq c} \frac{\bm{R}_{cc'}}{R_{cc'}^3}\,Z_{c'} .
\end{equation}
Aquí $\alpha_c$ es el tensor de polarizabilidad del centro, $\bm{\mu}_c^{\circ}$ su momento dipolar estático, y $C(r,\rho)$ una función de corte que regulariza el campo a distancias cortas.

\subsection{Imposición de la simetría atómica}

Tres condiciones reducen \eqref{eq:cpp_general} al caso que nos interesa.

En primer lugar, un átomo aislado posee un solo centro polarizable, de modo que la suma externa sobre $c$ desaparece y el segundo sumando de \eqref{eq:cpp_campo}, que representa el campo generado por los demás núcleos, se anula de manera idéntica. El vector $\bm{r}_{ci}$ se reduce al radio vector del electrón medido desde el núcleo, y el campo efectivo queda
\begin{equation}
    \label{eq:cpp_campo_atomico}
    \bm{f} = \sum_i \frac{\bm{r}_i}{r_i^3}\,W(r_i),
\end{equation}
donde $W(r)$ es la función de corte del campo, que se anula en el interior del core y tiende a la unidad fuera de él.

En segundo lugar, el core de nuestros sistemas es una capa cerrada y por tanto esféricamente simétrico, según el teorema de Unsöld de la Sección \ref{sec:unsold}. Una distribución de carga esférica no admite momento dipolar permanente, de manera que $\bm{\mu}^{\circ} = \bm{0}$ y los dos últimos términos de \eqref{eq:cpp_general} se desvanecen. Esta es la diferencia esencial con el caso molecular, donde los enlaces deforman la nube electrónica de forma permanente.

En tercer lugar, la isotropía de un core esférico reduce el tensor de polarizabilidad a un escalar, $\alpha_c \to \alpha_d$. Con estas tres condiciones,
\begin{equation}
    \label{eq:cpp_colapsado}
    V_{\text{pol}} = -\frac{1}{2}\alpha_d \,(\bm{f}\cdot\bm{f}).
\end{equation}

\subsection{Los términos de uno y de dos cuerpos}

Sustituir \eqref{eq:cpp_campo_atomico} en \eqref{eq:cpp_colapsado} produce un producto de dos sumatorias que conviene separar según los índices coincidan o no,
\begin{equation}
    V_{\text{pol}} = -\frac{1}{2}\alpha_d \left( \sum_i \frac{\bm{r}_i}{r_i^3}W(r_i) \right)\cdot\left( \sum_j \frac{\bm{r}_j}{r_j^3}W(r_j) \right).
\end{equation}

Los términos diagonales, con $i=j$, involucran el producto de un vector consigo mismo. Usando $\bm{r}_i\cdot\bm{r}_i = r_i^2$, la potencia se simplifica a $r_i^2/r_i^6 = 1/r_i^4$ y queda un operador estrictamente de un cuerpo,
\begin{equation}
    \label{eq:cpp_un_cuerpo}
    V_{\text{pol}}^{(1)} = -\frac{1}{2}\alpha_d \sum_i \frac{W^2(r_i)}{r_i^4},
\end{equation}
En nuestra implementación el cuadrado $W^2(r)$ se toma como $W_6(r/r_c) = 1 - e^{-(r/r_c)^6}$, de modo que \eqref{eq:cpp_un_cuerpo} coincide con la forma \eqref{eq:vpol_cutoff} del Capítulo \ref{ch:correlacion-estructura-fina} y con el potencial central que se suma a la atracción nuclear en la Ecuación \eqref{eq:vpol_efectivo}, que es el que interviene en el ciclo SCF. Con esa elección el corte del campo es $W = W_6^{1/2}$, y es este el que aparecería en el término de dos cuerpos. La dependencia $r^{-4}$, introducida en el Capítulo \ref{ch:correlacion-estructura-fina} por vía electrostática, emerge aquí como consecuencia puramente algebraica del producto punto.

Los términos cruzados, con $i \neq j$, sobreviven en cambio como un operador de dos cuerpos. Para la capa de valencia $np^2$ solo hay dos electrones y por tanto dos combinaciones cruzadas, $(1,2)$ y $(2,1)$, que al ser el producto punto conmutativo se suman y cancelan el factor $1/2$ frontal:
\begin{equation}
    \label{eq:cpp_dos_cuerpos}
    V_{\text{pol}}^{(2)} = -\alpha_d\,\frac{\bm{r}_1 \cdot \bm{r}_2}{r_1^3 r_2^3}\,W(r_1)W(r_2).
\end{equation}

Este término tiene una interpretación física propia y no es un residuo despreciable: describe el apantallamiento dieléctrico que el core interpone entre los dos electrones de valencia, ya que el dipolo inducido por uno de ellos es sentido por el otro. Se suma al operador de repulsión $1/r_{12}$ y, al depender del ángulo entre ambos radios vectores, contribuye a la parte anisótropa de la interacción, reduciendo la separación entre los términos del multiplete. Su omisión en un tratamiento que retenga únicamente \eqref{eq:cpp_un_cuerpo} constituye, por tanto, una aproximación que conviene tener presente al interpretar los desdoblamientos calculados.

\section{Costo Computacional del Cálculo CI}
\label{sec:costo_ci}

Las curvas de convergencia de la Sección \ref{sec:convergencia_ci} exigen repetir el cálculo de interacción de configuraciones para tamaños crecientes del espacio activo, y conviene documentar qué recurso limita ese barrido. El costo no lo domina la diagonalización, cuyo espacio es de unas pocas miles de configuraciones, sino el ensamblado previo de las integrales de Slater $R^k$ entre orbitales virtuales. Su número crece como la cuarta potencia del tamaño de la base de virtuales, de modo que es este ensamblado el que fija tanto el tiempo de pared como la huella de memoria.

\begin{table}[htbp]
    \centering
    % GENERADO por benchmarks/np2_sequence/tesis/etapa3_tablas.jl; no editar a mano.
% Procedencia:
%   CI carbon_rohf_results_3P_R30.0.jld2 m=4 (sha256 e8c5884b, commit 2479092)
%   CI carbon_rohf_results_3P_R30.0.jld2 m=8 (sha256 e8c5884b, commit 2479092)
%   CI carbon_rohf_results_3P_R30.0.jld2 m=12 (sha256 e8c5884b, commit 2479092)
%   CI carbon_rohf_results_3P_R30.0.jld2 m=16 (sha256 e8c5884b, commit 2479092)
%   CI carbon_rohf_results_3P_R30.0.jld2 m=20 (sha256 e8c5884b, commit 2479092)
%   CI carbon_rohf_results_3P_R30.0.jld2 m=24 (sha256 e8c5884b, commit 6dfb84b)
%   CI carbon_rohf_results_3P_R30.0.jld2 m=28 (sha256 e8c5884b, commit 6dfb84b)
%   CI carbon_rohf_results_3P_R30.0.jld2 m=32 (sha256 e8c5884b, commit 6dfb84b)
%   CI carbon_rohf_results_3P_R30.0.jld2 m=36 (sha256 e8c5884b, commit 6dfb84b)
%   CI carbon_rohf_results_3P_R30.0.jld2 m=40 (sha256 e8c5884b, commit 6dfb84b)
%   CI silicon_rohf_results_3P_R30.0.jld2 m=4 (sha256 09adeb32, commit 2479092)
%   CI silicon_rohf_results_3P_R30.0.jld2 m=8 (sha256 09adeb32, commit 2479092)
%   CI silicon_rohf_results_3P_R30.0.jld2 m=12 (sha256 09adeb32, commit 2479092)
%   CI silicon_rohf_results_3P_R30.0.jld2 m=16 (sha256 09adeb32, commit 2479092)
%   CI silicon_rohf_results_3P_R30.0.jld2 m=20 (sha256 09adeb32, commit 2479092)
%   CI germanium_rohf_results_3P_R30.0.jld2 m=4 (sha256 7567c491, commit 2479092)
%   CI germanium_rohf_results_3P_R30.0.jld2 m=8 (sha256 7567c491, commit 2479092)
%   CI germanium_rohf_results_3P_R30.0.jld2 m=12 (sha256 7567c491, commit 2479092)
%   CI germanium_rohf_results_3P_R30.0.jld2 m=16 (sha256 7567c491, commit 2479092)
%   CI germanium_rohf_results_3P_R30.0.jld2 m=20 (sha256 7567c491, commit 2479092)
%   CI tin_rohf_results_3P_R30.0.jld2 m=4 (sha256 7db46b8f, commit 2479092)
%   CI tin_rohf_results_3P_R30.0.jld2 m=8 (sha256 7db46b8f, commit 2479092)
%   CI tin_rohf_results_3P_R30.0.jld2 m=12 (sha256 7db46b8f, commit 2479092)
%   CI tin_rohf_results_3P_R30.0.jld2 m=16 (sha256 7db46b8f, commit 2479092)
%   CI tin_rohf_results_3P_R30.0.jld2 m=20 (sha256 7db46b8f, commit 2479092)
%   lmax = 3. t es incremental (cada tamano reutiliza las R^k del anterior) salvo donde resultados_ci.toml marca cache_heredada = false: primer tamano o corrida reanudada.
\begin{tabular}{rrrrrrrrr}
    \toprule
    $m$ & $R^k$ C & $R^k$ Si & $R^k$ Ge & $R^k$ Sn & $t$ C (s) & $t$ Si (s) & $t$ Ge (s) & $t$ Sn (s) \\
    \midrule
    4 & 4548 & 4738 & 5058 & 5378 & 4 & 0 & 4 & 10 \\
    8 & 65718 & 66402 & 67554 & 68706 & 4 & 3 & 15 & 33 \\
    12 & 324056 & 325538 & 328034 & 330530 & 13 & 12 & 59 & 112 \\
    16 & 1012266 & 1014850 & 1019202 & 1023554 & 33 & 32 & 156 & 285 \\
    20 & 2455212 & 2459202 & 2465922 & 2472642 & 70 & 68 & 316 & 598 \\
    24 & 5069918 & -- & -- & -- & 273 & -- & -- & -- \\
    28 & 9365568 & -- & -- & -- & 218 & -- & -- & -- \\
    32 & 15943506 & -- & -- & -- & 369 & -- & -- & -- \\
    36 & 25497236 & -- & -- & -- & 586 & -- & -- & -- \\
    40 & 38812422 & -- & -- & -- & 869 & -- & -- & -- \\
    \bottomrule
\end{tabular}

    \caption{Número de integrales $R^k$ almacenadas en caché y tiempo de pared por tamaño del espacio activo, para los cuatro elementos de la secuencia. Los tiempos son incrementales: cada tamaño reutiliza las integrales calculadas para el anterior, salvo en el primer punto de cada curva y en las corridas reanudadas.}
    \label{tab:convergencia_ci_costo}
\end{table}

El Cuadro \ref{tab:convergencia_ci_costo} cuantifica ambos recursos. El conteo de integrales es prácticamente el mismo para los cuatro elementos, ya que depende del tamaño de la base y no de la carga nuclear, y pasa de unos cuatro mil quinientos valores en $m=4$ a treinta y ocho millones en $m=40$. Con un ocupamiento medido de unos diecisiete bytes por casilla, el caso extremo del carbono requiere del orden de $1.1$ gigabytes de memoria residente, que llegan a $1.7$ durante los redimensionamientos de la tabla asociativa. El tiempo de pared del mismo punto es de unos quince minutos, y la curva completa del carbono hasta $m=40$ tomó cerca de treinta y nueve minutos.

La comparación entre elementos muestra que el tiempo sí depende de $Z$, aunque el número de integrales no. A $m=20$ el silicio emplea sesenta y ocho segundos y el estaño quinientos noventa y ocho, casi nueve veces más, pese a evaluar una cantidad de integrales que difiere en menos del uno por ciento. La razón es el número de puntos de la malla radial, que debe crecer con $Z$ para resolver la contracción de las capas internas: cada integral $R^k$ se evalúa sobre esa malla, de modo que el costo por integral escala con la discretización y no con el espacio activo. Esta separación entre los dos factores es la que permite reutilizar la caché entre tamaños consecutivos y hace viable el barrido completo.

\section{Arquitectura Computacional y Repositorio de Código}
\label{sec:repositorio}

Para salvaguardar la pureza teórica del manuscrito y mantener el rigor físico, se ha omitido la inclusión directa de código fuente a lo largo de esta tesis. Sin embargo, todos los métodos, derivaciones y algoritmos descritos (generación de mallas nodales, ortogonalización de Gram-Schmidt, rutinas de Campo Autoconsistente y resolución de la ecuación de Poisson por Galerkin) fueron programados y ensamblados desde cero en una librería de desarrollo propio bautizada como \textbf{AtomicSplines}.

El motor está escrito en \textbf{Julia}, un lenguaje compilado de cómputo científico. Inicialmente se exploró el uso de estructuras de datos esparsas especializadas (como las provistas por \texttt{BandedMatrices.jl}) debido a que el método de Galerkin con B-splines produce naturalmente matrices diagonales por bandas. No obstante, se determinó empíricamente que el sobrecosto de gestión de memoria (\textit{overhead}) de estas estructuras era prohibitivo para la escala de nuestras discretizaciones. Por lo tanto, se optó por una arquitectura mucho más veloz basada en el ecosistema básico de \texttt{LinearAlgebra.jl} auxiliado en ciertas rutinas críticas por \texttt{StaticArrays.jl}, lo que acelera drásticamente las operaciones algebraicas en memoria contigua.

Las diagonalizaciones del ciclo SCF y del cálculo de interacción de configuraciones son completas y densas, mediante las rutinas de LAPACK que expone la biblioteca estándar de álgebra lineal de Julia. Con bases de unos cientos de funciones y espacios de configuraciones de unos pocos miles, su costo no domina el cálculo, que como muestra la Sección \ref{sec:costo_ci} lo fija el ensamblado de las integrales $R^k$. Los símbolos $3j$ y $6j$ que requiere el cálculo de interacción de configuraciones se evalúan con la biblioteca \texttt{WignerSymbols.jl}.

Fomentando el paradigma de ciencia abierta y reproducibilidad, el código fuente completo, junto con los módulos de álgebra de tensores y la documentación técnica, se encuentra versionado y disponible públicamente en el siguiente repositorio de GitHub:

\begin{center}
    \url{https://github.com/recore799/AtomicSplines.git}
\end{center}
